\documentclass[11pt]{article}

\usepackage[margin=1in]{geometry}
\usepackage{xcolor}
\usepackage{graphicx}
\usepackage{booktabs}
\usepackage{amsmath}
\usepackage{enumitem}
\usepackage{listings}
\usepackage[hidelinks]{hyperref}

\newcommand{\file}[1]{\texttt{\detokenize{#1}}}
\newcommand{\code}[1]{\texttt{\detokenize{#1}}}

\newsavebox{\gpbox}
\newenvironment{finding}[1]{%
  \begin{lrbox}{\gpbox}\begin{minipage}{0.94\linewidth}%
  \textbf{\textcolor{green!35!black}{#1}}\par\smallskip\small}%
 {\end{minipage}\end{lrbox}%
  \begin{center}\setlength{\fboxrule}{0.7pt}%
  \fcolorbox{green!35!black}{green!5}{\usebox{\gpbox}}\end{center}}
\newenvironment{caveat}[1]{%
  \begin{lrbox}{\gpbox}\begin{minipage}{0.94\linewidth}%
  \textbf{\textcolor{red!55!black}{#1}}\par\smallskip\small}%
 {\end{minipage}\end{lrbox}%
  \begin{center}\setlength{\fboxrule}{0.7pt}%
  \fcolorbox{red!55!black}{red!4}{\usebox{\gpbox}}\end{center}}

\title{\vspace{-1.5em}Hydration-site clustering: old vs new\\[0.3em]
\large The same CRY1/AN139 trajectory analysed two ways}
\author{}
\date{}

\begin{document}
\maketitle
\vspace{-2.5em}

\section*{Summary}

The original GRAND clustering and the pipeline's \file{ev71_postprocess.py} were run on
\textbf{the identical 2500-frame trajectory}, with the same sphere, cutoff and frame
sampling. Only the analysis differs, so any difference is method, not sampling.

\textbf{They agree on every persistent hydration site.} Above an occupancy of 0.4 the two
site sets are identical. Where they overlap, sites coincide to \textbf{0.36 \AA{}}.

\textbf{The old method reports 302 clusters, the new one 176.} The 126-cluster difference is
almost entirely one-frame noise: the extra clusters have median occupancy \textbf{0.00}.
This is the ghost-exclusion fix removing spurious water observations, and it matters because
those clusters are what depressed the precision and Tanimoto figures in the original study.

\section{Setup}

\begin{center}\small
\begin{tabular}{@{}l l@{}}
\toprule
Trajectory & \file{GCMC_Grand/CRY1AN139-raw.dcd}, 2500 frames (complete production) \\
Old result & \file{CRY1AN139-lig-clusts.pdb}, already present, from GRAND \code{cluster_waters} \\
New result & \file{ev71_postprocess.py}, \code{--cluster-stride 1} (full parity) \\
Common settings & 10 \AA{} sphere on the ligand, 2.4 \AA{} average-linkage cutoff \\
\midrule
New-method cost & 71\,050 water observations, 20.2 GB distance matrix, 16 min \\
\bottomrule
\end{tabular}
\end{center}

The one substantive difference between the implementations is ghost handling. GRAND shifts
inactive ghost waters \emph{before} imaging, which can bring non-interacting molecules back
toward the ligand; the new script images first, shifts after, and additionally excludes
inactive ghosts frame by frame.

\section{Result 1 --- the site sets}

\begin{center}
\includegraphics[width=0.98\linewidth]{AN139method_occupancy.png}
\end{center}

Left: every cluster, ranked by occupancy. The two curves lie on top of one another for the
first $\sim$130 sites and then diverge --- the old method continues with a long tail of
clusters that the new one does not produce. Right: the same thing binned (note the log
scale); the excess is entirely in the lowest bin.

\begin{center}\small
\begin{tabular}{@{}l r r r@{}}
\toprule
& all clusters & occupancy $\ge 0.2$ & median occupancy of the extras \\
\midrule
old (GRAND) & 302 & 49 & \\
new (pipeline) & 176 & 44 & \\
\midrule
matched & 168 & 39 & \\
old only & 134 & & \textbf{0.00} \\
new only & 8 & & 0.08 \\
\bottomrule
\end{tabular}
\end{center}

The 134 old-only clusters have a median occupancy of \textbf{0.00} --- they are occupied in
about one frame out of 2500. They are not distant artifacts: their distance to the ligand
(median 3.8 \AA{}) is the same as the matched sites'. They are one-off water positions inside
the sphere that the ghost-exclusion step removes.

\section{Result 2 --- how well they agree}

\begin{center}
\includegraphics[width=0.98\linewidth]{AN139method_agreement.png}
\end{center}

Left: the 39 matched persistent sites, mean separation \textbf{0.36 \AA{}}. For scale, two
independent GCMC \emph{simulations} of this system match at $\sim$1.1--1.3 \AA{}. Two
analyses of the \emph{same} trajectory should agree far better than that, and they do.

Right: sweeping the occupancy cut. This is the most informative panel.

\begin{center}\small
\begin{tabular}{@{}r r r r r r r@{}}
\toprule
cut & $n_{\text{old}}$ & $n_{\text{new}}$ & matched & recall & precision & Tanimoto \\
\midrule
0.00 & 302 & 176 & 168 & 0.556 & 0.955 & 0.542 \\
0.10 & 133 & 125 & 122 & 0.917 & 0.976 & 0.897 \\
0.20 & 49 & 44 & 39 & 0.796 & 0.886 & 0.722 \\
0.30 & 11 & 18 & 11 & 1.000 & 0.611 & 0.611 \\
\textbf{0.40} & \textbf{9} & \textbf{9} & \textbf{9} & \textbf{1.000} & \textbf{1.000} & \textbf{1.000} \\
0.50 & 4 & 4 & 4 & 1.000 & 1.000 & 1.000 \\
\bottomrule
\end{tabular}
\end{center}

\begin{finding}{Two things to take from this table}
\textbf{Precision is 0.955 even with no occupancy cut.} Almost every site the new method
reports is also in the old set. The new method is not finding different water --- it is
declining to report spurious clusters. The disagreement is one-directional.

\textbf{Above a 0.4 cut the sets are identical.} Every persistent hydration site, the kind
that actually gets reported, is found by both. The entire discrepancy lives in clusters that
appear in a handful of frames.
\end{finding}

\section{Result 3 --- where the disagreement sits}

\begin{center}
\includegraphics[width=0.98\linewidth]{AN139method_disagreement.png}
\end{center}

Every old-only cluster (red) sits at the bottom of the occupancy axis, spread across the same
range of ligand distances as the matched sites. There is no spatial signature --- these are
not, for instance, ghosts piled up next to the ligand. They are simply water seen once.

\section{Why this matters for the published numbers}

The original study reported, for AN139 against crystallographic waters, a recall of 0.870
against a precision of \textbf{0.211} and a Tanimoto of \textbf{0.205}, and explained the low
values as follows:

\begin{quote}\small
``besides matched clusters there are many additional unmatched clusters, resulting in a low
value of Precision and Tanimoto similarity coefficient. This observation is expected, as
crystallographic structures provide a static and ensemble-averaged representation of
hydration and may not resolve transient or weakly occupied water molecules.''
\end{quote}

That reasoning is sound as far as it goes, but this comparison shows a second contribution
the study could not have seen: of the 302 clusters, \textbf{134 are one-frame artifacts of
ghost handling}, not transient physical hydration. Recomputing those metrics on the cleaned
176-cluster set would raise precision and Tanimoto without changing recall.

\begin{caveat}{What this does \emph{not} show}
Only one system and one trajectory. And this compares two \emph{analyses}, not two
\emph{simulation engines} --- it says nothing about whether GRAND and Loch sample the same
water, only that given identical frames they identify the same sites.
\end{caveat}

\section*{Update (2026-08-06): the engine comparison has since been done}

The matched-length production runs this document said did not yet exist have since been
produced (\file{an139\_loch\_full/}, 2 replicas $\times$ 2500 frames). Comparing the two
\emph{engines} --- same clustering on both sides, both site sets superposed into a common
frame by Kabsch on 476 C$\alpha$ atoms --- gives:

\begin{center}\small
\begin{tabular}{@{}lrrr@{}}
\toprule
AN139, occ $\ge$ 0.2 & matched & Tanimoto & mean $d$ \\
\midrule
GRAND vs Loch rep1 & 32 & \textbf{0.615} & 0.95 \AA{} \\
GRAND vs Loch rep2 & 28 & 0.528 & 0.85 \AA{} \\
Loch rep1 vs rep2 (replicate floor) & 28 & 0.571 & 0.58 \AA{} \\
\bottomrule
\end{tabular}
\end{center}

\noindent
Cross-engine agreement equals or exceeds the replicate floor: the two engines produce
\textbf{statistically indistinguishable} hydration-site networks.

One caution learned the hard way, directly relevant to this document's method. An earlier
version of that engine comparison reported Tanimoto 0.254 and concluded the engines disagreed
badly. That was wrong. It compared GRAND's \emph{old} clustering against Loch's \emph{new}
clustering (two variables at once), and --- dominant effect --- never reconciled the
coordinate frames. \file{ev71\_postprocess.py} superposes each trajectory onto \emph{its own}
frame 0, leaving a \textbf{3.5 \AA{} rigid-body offset} between independent runs, against a
2.0 \AA{} match cutoff. Fixing the clustering alone moved 0.254 $\rightarrow$ 0.273; fixing
the frame moved it to 0.615.

\textbf{Site sets from independent runs must be brought into a common frame before
comparison.} Artifacts: \file{an139\_fair\_engine\_comparison/}.

\section{Reproducing this}

\begin{lstlisting}[basicstyle=\ttfamily\footnotesize]
python ev71_postprocess.py \
  --topology  <GCMC_Grand>/CRY1AN139-ghosts.pdb \
  --trajectory <GCMC_Grand>/CRY1AN139-raw.dcd \
  --ghost-file <GCMC_Grand>/gcmc-ghost-wats.txt \
  --output-dir new --prefix AN139new \
  --sphere-radius 10.0 --cluster-cutoff 2.4 --cluster-stride 1 \
  --max-distance-memory-gb 64

python compare_water_sites.py \
  --sites old_grand <GCMC_Grand>/CRY1AN139-lig-clusts.pdb \
  --sites new_loch  new/AN139new-lig-clusts.pdb \
  --reference old_grand --plots \
  --ligand-reference <GCMC_Grand>/CRY1AN139-ghosts.pdb \
  --output-dir . --prefix AN139method
\end{lstlisting}

\noindent\small
\code{--cluster-stride 1} is needed for parity with GRAND and costs 20 GB and 16 minutes. A
stride of 5 gives the same persistent sites in 30 seconds and 0.6 GB; the stride is recorded
in the metrics JSON, so the approximation is never silent.

\end{document}
